\subsection{The exact coprime mean in the one-long-variable case}
\label{sec:one-long-coprime-mean}
Here is an explicit reconciliation needed when applying L23,
Lemma~3.4, in the one-long-variable alternatives above.
Write $z_0=x^{1/(\log\log x)^3}$ and let $R(n)$ be the indicator
that all prime factors of $n$ exceed $z_0$. The printed main
sum in that lemma is unrestricted, whereas our discrepancy uses
the exact coprime sum. For $q\le x^{3/5}$ and a long-variable
endpoint $t\gg x^{1-10\zeta}$, the difference is bounded by
\[
 \frac{1}{\varphi(q)}
 \sum_{\substack{p\mid q\\p>z_0}}(t/p+1)
 \ll \frac tq\left\{\frac{\log^C x}{z_0}
                         +\frac{\log^C x}{t}\right\}
 \ll_A\frac{t}{q\log^A x}
\]
for every fixed $A$. Indeed, a rough integer not coprime to
$q$ is divisible by one of the displayed primes; use the union
bound, $q/\varphi(q)\ll\log^C x$, and
$\omega(q)\ll\log x$. The bound applies at both endpoints
of any long-variable interval. Partial summation allows the
logarithm and Mellin twist already present in the proof.

For a fixed short product $m$, if $(m,q)>1$, the progression
and exact coprime terms both vanish, since $(N,q)=1$.
Otherwise the long-variable residue is $N\overline m\pmod q$
and is primitive. Summation over the short products and
moduli costs only powers of the logarithm, by the fixed
divisor moment estimates for $\sum_{m\asymp M}\tau(m)^C/m$
and $\sum_{q\le Q}\tau(q)^C/q$. Increasing $A$ therefore
absorbs this correction in the stated remainder. This
reconciliation changes neither a distribution level nor a
sieve constant.
